\chapter[Conclusion and Future Work]{Conclusion and Future Work}\label{chapter:runtime-improvement}\label{chapter:discussion}\label{chapter:conclusion}

\begin{quote}\itshape
The objective of this chapter is to consolidate the thesis and identify
directions for further research. The chapter begins by summarising the evidence
from the three experiments, relating it to the research questions, and
restating the contributions and evidence boundaries of the study. It then
develops future work from the limitations identified in the evaluation.
Finally, it closes the thesis by linking the present findings to the next
stages of validation, cross-robot comparison, and real-world extension.
\end{quote}

\section[Summary]{Summary}

This thesis investigated whether large language models can generate reusable
robot-specific software for high-level embodied agents. Auto-Adapter derives
capability interfaces from robot descriptions and documented task requirements,
then implements them as Python drivers. An independent Harness assesses driver
behaviour using MuJoCo state, while \mbox{ReCAP} uses these interfaces to execute
downstream tasks whose completion is judged separately. Through this framework,
the software connecting a high-level planner to a robot becomes an explicit
object of generation and behavioural evaluation.

For RQ1, the PiPER experiment showed that Auto-Adapter can generate robot
drivers for reuse across tasks. Using a supplied controller skeleton, two of
five independent runs produced drivers that passed their complete capability
test suites. The exported drivers were reused unchanged across reaching,
pushing and pick-and-place tasks, with reaching and pushing succeeding in all
evaluated layouts. With the driver unchanged, pick-and-place tasks passed
after the high-level controller model was replaced with the more capable
\mbox{GPT-6 Astra}. This shows that the driver provided the control capabilities needed
for the evaluated tasks, supporting its functional effectiveness and value for
cross-task reuse.

For RQ2, AutoAdapter-Bench compared driver synthesis with each robot's
capability interface and behavioural tests fixed across models and independent
runs. Models generated and revised only the driver implementation.
Downstream model performance was evaluated separately using a fixed validated
driver, with task conditions and the ReCAP planning procedure held constant.
The results showed a distinction between driver synthesis
and driver use. Haiku produced no driver that passed its complete test suite,
yet completed some tasks through supplied drivers. Validation feedback helped
correct implementation errors, but successive repairs did not consistently
improve behavioural performance. By combining validation outcomes, task-success
rates and resource use, the benchmark informed model selection for driver
synthesis and high-level control.

For RQ3, the Franka Panda, LEAP Hand and Skydio X2 cases showed uneven task
support across robot configurations. The Panda driver supported multiple
manipulation tasks, and the Skydio driver supported several flight tasks.
LEAP's accurate local motion did not translate into reliable object
manipulation. These cases exposed limitations in sustained contact,
coordinated motion and orientation control that individual capability tests
did not fully assess. Each case evaluated one driver, and some results were
reassessed under revised criteria, limiting conclusions about consistency
across robots.

Together, these studies contribute a framework for generating robot drivers
from task requirements and a benchmark that separates implementation quality
from downstream use. The experiments show that independent measurement must be
paired with tests that accurately express the intended behaviour. Generated
criteria therefore require human review, and downstream task success must be
assessed separately from capability conformance. This separation of
responsibilities provides a practical basis for generating, inspecting and
reusing robot software through Auto-Adapter. Performance on real hardware and
sim-to-real transfer still require independent validation.

\section[Future Work]{Future Work}

\textbf{First, behavioural validation should cover a broader range of
conditions.} Future studies could introduce unseen tasks, varied initial
states, scene changes, and repeated resets to test whether an accepted driver
remains reliable beyond its validation suite. Validation thresholds and
boundary cases should also be calibrated further, distinguishing stable
drivers from implementations that fit only the evaluated cases.

\textbf{Second, failures associated with perception, control structure, and
driver synthesis should be separated.} The present experiments mainly use
structured state observations supplied by the simulator. Future work could
compare structured state, noisy state, and RGB or RGB-D observations, while
varying control priors separately and holding the remaining conditions fixed.
This would provide clearer evidence about the source of each failure.

\textbf{Third, long-horizon planning and the role of Experience require
further investigation.} With the driver, tasks, and Harness held fixed,
experiments could test whether explicit subgoal records, feedback use, and
recovery mechanisms reduce false completion and improve task success. The
effect of Experience would require a matched no-Experience condition and
multiple independent runs. Loading Experience across runs establishes an
information-transfer mechanism, but does not by itself establish improved
performance.

\textbf{Fourth, cross-robot evaluation requires a more closely matched
design.} Future experiments should hold the model, task structure, resource
budget, and validation standard as constant as possible, making outcomes more
comparable across robot configurations. Any investigation of morphology would
also need to separate morphology from differences in control interfaces,
observations, and tasks.

\textbf{Finally, simulator evidence should be extended to real execution in
stages.} Future work should first connect the framework to genuine vendor
software development kits and robot-specific Translation Layers. Within that
extension, the Model Hardware Standard described in
Section~\ref{sec:interfaces-drivers} could be evaluated as a common interface
between the generated robot-specific driver and device-specific execution,
rather than as a replacement for the underlying vendor integration. A matched
comparison should hold the robot configuration, capability interface, tasks,
generation model and budget, and trusted Harness measurement definitions and
verdict rules fixed while comparing direct-SDK and MHS-backed execution routes.
It should measure integration effort, behavioural correctness, communication
latency, and recovery from device errors. Physical experiments would
additionally examine calibration, sensor noise, contact behaviour, and
violations of declared interface limits. The present thesis does not integrate
MHS or evaluate an MHS-backed execution route, so it provides no evidence that
this route improves interoperability or physical safety. Simulator validation
can serve as a filter before either route reaches real hardware, but it is not
sufficient evidence for physical deployment.
